\documentclass[11pt]{article}
\usepackage[T1]{fontenc}
\usepackage[utf8]{inputenc}
\usepackage{lmodern,amsmath,amssymb,amsthm,mathtools,booktabs,array,microtype}
\usepackage[margin=1.05in]{geometry}
\usepackage[colorlinks=true,linkcolor=blue,citecolor=blue,urlcolor=blue,breaklinks=true]{hyperref}
\usepackage{enumitem}
\setlist{itemsep=3pt,topsep=5pt}
\newcommand{\workbenchurl}{https://github.com/KokunoYumeto/algebraic-cycles-and-motives-workbench}
\newcommand{\recordurl}{https://github.com/KokunoYumeto/mathematics-workbenches}
\newcommand{\CH}{\operatorname{CH}}
\newcommand{\Alb}{\operatorname{Alb}}
\newcommand{\Pic}{\operatorname{Pic}}
\newcommand{\NS}{\operatorname{NS}}
\newcommand{\Hom}{\operatorname{Hom}}
\newcommand{\cl}{\operatorname{cl}}
\newcommand{\Gr}{\operatorname{Griff}}
\newcommand{\Z}{\mathbb Z}
\newcommand{\Q}{\mathbb Q}
\newcommand{\C}{\mathbb C}
\newcommand{\F}{\mathbb F}
\newcommand{\PP}{\mathbb P}
\newcommand{\alg}{\mathrm{alg}}
\newcommand{\et}{\mathrm{\acute et}}
\DeclareMathOperator{\rad}{rad}
\DeclareMathOperator{\id}{id}
\newtheorem{theorem}{Theorem}[section]
\newtheorem{proposition}[theorem]{Proposition}
\newtheorem{lemma}[theorem]{Lemma}
\newtheorem{corollary}[theorem]{Corollary}
\theoremstyle{definition}
\newtheorem{definition}[theorem]{Definition}
\newtheorem{example}[theorem]{Example}
\theoremstyle{remark}
\newtheorem{remark}[theorem]{Remark}
\title{Grothendieck's Albanese equivalence and his conjectures on algebraic cycles}
\author{}
\date{23 November 2026}
\begin{document}
\maketitle
\begin{abstract}
Grothendieck's Montpellier notes propose a cohomological consequence of Albanese equivalence and a Lefschetz theory for cycles modulo algebraic equivalence. We prove the first implication in every positive characteristic, with an integral hypothesis, by specializing the entire Albanese-kernel relation and applying Bloch's prime-to-characteristic torsion theorem. An elliptic curve shows why the hypothesis cannot be rationalized. Over the complex numbers, Clemens's theorem and an iterated blow-up construction give the exact universal range of the proposed Lefschetz surjectivity; the Hodge--Riemann relations identify the obstruction to primitive positivity. We also obtain bounds for integral scalar factorizations through products of Jacobians, using the work of Oort--Ueno, Beauville, Alexeev--Donagi--Farkas--Izadi--Ortega, and Scavia--Suzuki. The remaining surface and threefold obstructions are confined respectively to the primes two, and two and three. The mathematics and text were produced by GPT-6 Astra in ChatGPT Pro mode, with separate checks by the same model as described at the end; no human has checked them.
\end{abstract}

\section{Grothendieck's questions and the results}

Grothendieck considered two different ways of extracting linear information from algebraic cycles. Quotienting by algebraic equivalence retains more information than cohomology, whereas quotienting algebraically trivial cycles by relations coming from Albanese kernels is intended to isolate information of abelian type. These distinctions matter integrally. The results below describe where the proposed comparisons hold and where they fail.

The archive locators need to be distinguished. In cote~14, \emph{Th\'eorie des motifs II}, the handwritten Theorem~3.1 is on sheet~37 and its argument continues on sheet~38. The displayed assertion is
\begin{equation}\label{eq:archive}
             u_Z=0\quad\Longrightarrow\quad a(Z)=0.
\end{equation}
The separate letter to Serre dated 27 August 1965 begins on sheet~43. There Grothendieck's Conjecture~A asks that
\[
 C^i(X)\xrightarrow{\;\xi^{n-2i}\;}C^{n-i}(X)
 \quad\text{``induit un isomorphisme''},
\]
where $C^i$ means cycles modulo algebraic equivalence with rational coefficients. Conjecture~B on sheet~44 asks for positive definiteness of the primitive middle intersection form; sheet~45 gives the proposed finite-generation consequences. These quotations and sheet numbers refer to the archive's printed numbering, not to the PDF page counter \cite{Grothendieck}.

\subsection{An integral implication for Albanese equivalence}

All Chow groups have integral coefficients unless a subscript $\Q$ is present. Let $X$ be smooth, projective and connected over an algebraically closed field $k$. Write $\CH^i(X)_{\alg}$ for algebraically trivial classes. Define $N^i(X)$ to be the subgroup generated by
\begin{equation}\label{eq:N}
 W_*\nu,\qquad W\in\CH^i(T\times X),\quad
 \nu\in\CH_0(T)^0,\quad \operatorname{alb}_T(\nu)=0,
\end{equation}
where $T$ ranges over smooth projective connected varieties. Set
\[
                  P^i(X)=\CH^i(X)_{\alg}/N^i(X).
\]
This is the zero-cycle-generated form of Albanese equivalence used here; no identification with the intersection of all regular-homomorphism kernels is implicit.

Let $A$ be an abelian variety. An integral correspondence $Z\in\CH^i(A\times X)$ gives a homomorphism
\[
 u_Z:A(k)\longrightarrow P^i(X),\qquad
 a\longmapsto [Z_a-Z_0].
\]
Indeed, the zero-cycle $[a+b]-[a]-[b]+[0]$ has zero Albanese class on $A$, so its image lies in $N^i(X)$. Subtracting the constant correspondence normalizes $Z$ without changing the relevant cohomological map.

\begin{theorem}[Albanese equivalence in positive characteristic]\label{thm:alb}
Let $k$ be algebraically closed of characteristic $p>0$, let $X/k$ be smooth projective connected of dimension $n$, let $A/k$ be an abelian variety, and let $1\leq i\leq n$. If $Z\in\CH^i(A\times X)$ is integral and $u_Z=0$ as a homomorphism of integral groups, then, for every prime $\ell\ne p$, its transpose induces the zero map
\[
 {}^tZ_*:H^{2(n-i)+1}_{\et}(X,\Q_\ell(n-i))
                  \longrightarrow H^1_{\et}(A,\Q_\ell).
\]
The corresponding map with $\Z_\ell$ coefficients is also zero. The conclusion need not hold if the premise is weakened to $u_Z\otimes\Q=0$.
\end{theorem}

The main issue is not the usual compatibility of correspondences with the torsion cycle map. A torsion class in the image of an Albanese kernel need not come from a torsion class in that kernel: a homomorphism $\Q\to\Q/\Z$ already illustrates the obstruction. We instead spread a finite expression \eqref{eq:N}, all its rational-equivalence witnesses, and the equation killing the resulting class. On a finite-field fibre every degree-zero zero-cycle is torsion, and Bloch's theorem separates its $\ell$-primary part. This proves a vanishing statement for the torsion cycle map before taking Tate modules.

Bloch proved the required correspondence, specialization and Albanese compatibilities in Propositions~3.5, 3.8 and 3.9 of \cite[pp.~115, 117--119]{Bloch}; his Theorem~4.2, p.~119, proves the prime-to-characteristic torsion theorem. The proof below uses precisely those results. It does not use the stronger theorem announced in his Theorem~4.1 as though that announcement were a proof.

\subsection{The complex Lefschetz assertions}

For a smooth projective complex $n$-fold put
\[
 A^r(X)=(\CH^r(X)/\CH^r(X)_{\alg})\otimes\Q,
 \qquad G^r(X)=\ker\bigl(A^r(X)\xrightarrow{\cl}H^{2r}(X,\Q(r))\bigr).
\]
Let $L$ denote intersection with an ample divisor class.

\begin{theorem}[Universal range modulo algebraic equivalence]\label{thm:range}
The following assertions hold.
\begin{enumerate}[label=(\roman*)]
\item For $0\leq i\leq n/2$, the map
$L^{n-2i}:A^i(X)\to A^{n-i}(X)$ is surjective for every smooth projective complex $n$-fold and every ample $L$ if and only if $i=0$ or $2i=n$. In those two cases it is an isomorphism. For each $1\leq i<n/2$ there is an $X$ for which the source is finite-dimensional and the target infinite-dimensional, for every ample $L$.
\item For $0\leq r\leq n$, each of the following properties holds for every such $n$-fold in codimension $r$ if and only if $r\in\{0,1,n\}$: injectivity of the rational cycle class map; coincidence of algebraic and numerical equivalence after tensoring with $\Q$; finite-dimensionality of $A^r(X)$; finite generation of $(\CH^r(X)/\CH^r(X)_{\alg})/\mathrm{torsion}$.
\item If $n=2m$ and $P^m_L(X)=\ker(L:A^m(X)\to A^{m+1}(X))$, then
\[
 q(x,y)=(-1)^m\deg(xy)
\]
is positive semidefinite on $P^m_L(X)$ and
\[
 \rad(q)=\ker\bigl(L:G^m(X)\to G^{m+1}(X)\bigr).
\]
Thus positive definiteness is exactly the injectivity of this last map. It holds universally for $m=0,1$. No assertion for all $m\geq2$ follows from parts~(i) and~(ii).
\end{enumerate}
\end{theorem}

The input for the negative assertions is Clemens's theorem on a very general quintic threefold \cite[Theorem~(0.2), equation~(0.1), p.~19]{Clemens}. Products with projective space settle part~(ii), but do not by themselves give a finite-dimensional source in every case of part~(i). A divisor incidence construction makes it possible to move the single infinite-dimensional Griffiths group upwards one degree at a time. Part~(iii) is a direct application of the Hodge--Riemann bilinear relations. It separates the positivity question from the failure of Lefschetz surjectivity.

For comparison, homological and numerical equivalence over $\C$ are known to agree in codimensions $0,1,2,n-1,n$, hence in all codimensions on varieties of dimension at most four. Kleiman records the divisor and curve cases and the dimension-at-most-four consequence \cite[Remarks~3.4 and~3.10, pp.~379--380]{Kleiman}; the codimension-two argument in arbitrary dimension is written out in Section~3.2. This is a statement about homological equivalence, and does not restore algebraic equivalence in Theorem~\ref{thm:range}.

\subsection{Integral factors of products of Jacobians}

For an abelian variety $A/k$ over an algebraically closed field, let
\[
 I_J(A)=\{a\in\Z:[a]_A=vu\text{ for homomorphisms }
 A\xrightarrow{u}\prod_jJ(C_j)\xrightarrow{v}A\},
\]
with smooth projective connected curves $C_j$. This is an ideal of $\Z$. The theorem below shows it is nonzero, so write $I_J(A)=s_J(A)\Z$ with $s_J(A)>0$. Define
\[
                    \delta(A)=\gcd\{\chi(A,L):L\in\Pic(A)\}.
\]
The invariant $s_J(A)$ is one exactly when $A$ is a direct factor, as an unpolarized abelian variety, of a finite product of Jacobians.

\begin{theorem}[Bounds for the Jacobian factor obstruction]\label{thm:scalar}
Over every algebraically closed field, in every characteristic, an abelian variety of dimension $g>0$ satisfies
\[
                     s_J(A)\mid(2g-1)!.
\]
If $g\leq3$, it also satisfies $s_J(A)\mid\delta(A)$. More precisely,
\[
 \begin{array}{c|c}
  \dim A&\text{bound}\\\hline
  2&s_J(A)\mid\gcd(2,\delta(A))\\
  3&s_J(A)\mid\gcd(6,\delta(A)).
 \end{array}
\]
Consequently an abelian surface with an ample line bundle of odd Euler characteristic is a direct factor of a product of Jacobians. The same holds for a threefold if its ample Euler characteristics have greatest common divisor prime to six.
\end{theorem}

These bounds are quantitative consequences of established constructions. The underlying integral factorization principle is already present in Voisin \cite[Proposition~2.1 and Corollary~2.5]{Voisin}, including the use of the Poincar\'e line bundle on $A\times\widehat A$. Beckmann--de Gaay Fortman \cite[Proposition~4.6 and Theorem~5.3]{Beckmann} give the related prime-to-prime isogeny criterion and bounds for the integral Hodge conjecture up to a factor. The exact scalar in the correspondence construction is the Matsusaka calculation in Scavia--Suzuki \cite[Proposition~3.4, equations~(3.1)--(3.3)]{ScaviaSuzuki}. Oort--Ueno's theorem handles principally polarized varieties of dimension at most three \cite[Theorem~4, p.~2 of the preprint]{OortUeno}. Beauville's generalized Pryms supply the factor two in dimensions four and five \cite[Proposition~2.2, Theorems~5.4 and~6.5(i)]{Beauville}; Alexeev--Donagi--Farkas--Izadi--Ortega supply the factor six in dimension six \cite[Corollary~2, p.~165]{ADFIO}. The deduction for nonprincipal polarizations uses $A\times\widehat A$ and the Euler characteristic of its Poincar\'e line bundle. In positive characteristic the curve cycles are specialized integrally.

The higher-dimensional question has changed substantially: Engel, de Gaay Fortman and Schreieder's preprint \cite[Theorem~1.1 and Corollary~1.5]{Engel} proves that the minimal curve class on a very general principally polarized complex abelian variety of dimension at least four is not algebraic. Together with the criterion in \cite{ScaviaSuzuki}, this obstructs integral Jacobian factors in those dimensions. Their subsequent preprint \cite[Theorem~1.2]{EngelSix} shows that six is the exact minimal algebraic multiple on a very general principally polarized sixfold. This result concerns a general sixfold, whereas $A\times\widehat A$ lies in a special locus; the low-dimensional bounds here retain their remaining factors two and three.

Recent preprints of OpenAI state the rational Hodge conjecture for CM abelian varieties \cite[Theorem~1.1, p.~1]{OpenAICM}, its Tate and Hodge-standard consequences for abelian varieties \cite[Corollaries~8.3--8.4, pp.~49--50]{OpenAICM}, and the rationality of Hodge-class pairings after good reduction \cite[Theorem~1.1, p.~1]{OpenAIRationality}. These statements concern rational cohomology classes and positivity modulo numerical equivalence. They do not determine the kernel modulo algebraic equivalence in Theorem~\ref{thm:range}, or the integral scalar ideal in Theorem~\ref{thm:scalar}. They are neighboring preprint claims, not inputs to the proofs here.

\section{Specializing an entire Albanese-kernel relation}

This section proves the torsion statement needed for Theorem~\ref{thm:alb}. For $\ell\ne\operatorname{char}k$, let
\[
 \lambda_X^i:\CH^i(X)[\ell^\infty]
      \longrightarrow H^{2i-1}_{\et}(X,\Q_\ell/\Z_\ell(i))
\]
be Bloch's torsion cycle map.

\begin{lemma}\label{lem:torsion}
If $k$ is algebraically closed of positive characteristic, then
\[
 \lambda_X^i\bigl(N^i(X)\cap\CH^i(X)[\ell^\infty]\bigr)=0.
\]
\end{lemma}
\begin{proof}
Choose $z$ in the displayed intersection and choose $r\geq0$ with $\ell^rz=0$. By definition there is a finite expression
\begin{equation}\label{eq:whole}
 z=\sum_{j=1}^a(W_j)_*\nu_j,\qquad
 \deg\nu_j=0,\qquad \operatorname{alb}_{T_j}(\nu_j)=0.
\end{equation}
The individual $\nu_j$ are not assumed torsion. Fix cycle representatives and rational-equivalence witnesses for both \eqref{eq:whole} and $\ell^rz=0$. A rational equivalence is a finite sum of divisors of rational functions on integral subvarieties; all the required witnesses therefore involve finitely many coefficients.

Let $B_j=\Alb(T_j)$ and fix Albanese maps $b_j:T_j\to B_j$ using base points. Descend $X,T_j,B_j,b_j,W_j,\nu_j,z$, all the witnesses, and the equalities $b_j(\nu_j):=\sum_x n_xb_j(x)=0$ in the group laws of the abelian varieties to a finitely generated subfield of $k$. Shrinking a smooth integral model $S$ of that field over a finite field, we obtain smooth projective families for $X$ and all $T_j$, abelian schemes for all $B_j$, and extensions of the maps and cycle data. We may shrink again so that the supports and rational functions defining the witnesses specialize properly. Thus the two Chow equalities and all the abelian-variety equalities hold on every remaining fibre. This assertion is about the chosen finite witnesses, not about specialization of an arbitrary infinite subgroup.

Choose a closed point $s\in S$ and pass to a geometric residue field $\bar\F_p$. One may realize the specialization through successive valuation rings associated with a regular flag through $s$, and then take algebraic closures of the fraction and residue fields. The constructions and equalities are compatible with each such base change. Denote special fibres by a subscript $s$; then
\begin{equation}\label{eq:specialwhole}
 z_s=\sum_j(W_{j,s})_*\nu_{j,s},\qquad
 \ell^rz_s=0,\qquad b_{j,s}(\nu_{j,s})=0.
\end{equation}

We must compare $B_{j,s}$ with the Albanese variety of $T_{j,s}$; equality of these two abelian schemes has not been assumed. On the geometric generic fibre the Albanese map induces an integral isomorphism
\[
 b_j^*:H^1(B_j,\Z_\ell)\xrightarrow{\sim}H^1(T_j,\Z_\ell).
\]
Smooth proper base change identifies these groups with the fibres of lisse $\Z_\ell$ sheaves, and the map between those sheaves is an isomorphism if it is one at a geometric point of the connected base. It follows that $b_{j,s}^*$ is also an integral isomorphism. By the universal property of the Albanese, $b_{j,s}$ factors through a homomorphism
\[
            \beta_j:\Alb(T_{j,s})\longrightarrow B_{j,s}.
\]
Its map on integral $\ell$-adic Tate modules is an isomorphism. Hence $\beta_j$ is an isogeny of degree prime to $\ell$: the rational isomorphism gives an isogeny and the integral isomorphism excludes an $\ell$-primary part of its kernel.

Every class in $\CH_0(T_{j,s})^0$ is torsion. Indeed, finitely many supporting points can be put on a smooth connected projective curve in $T_{j,s}$, and the class is then the pushforward of a degree-zero divisor. Such a divisor corresponds to a point of its Jacobian over $\bar\F_p$; that point is defined over a finite field and has finite order. (For a zero-dimensional $T_j$ the assertion is immediate.) Thus $\nu_{j,s}$ has a well-defined $\ell$-primary component. Equation \eqref{eq:specialwhole} and the prime-to-$\ell$ isogeny $\beta_j$ show that the Albanese class of this component is zero. Bloch's Theorem~4.2 \cite[p.~119]{Bloch} makes the Albanese map injective on prime-to-$p$ torsion. The $\ell$-primary component of $\nu_{j,s}$ is consequently zero.

It follows that each $\nu_{j,s}$ has order prime to $\ell$, and therefore so does the finite sum defining $z_s$. But \eqref{eq:specialwhole} also makes $z_s$ $\ell$-primary. Hence $z_s=0$ in $\CH^i(X_s)$.

Bloch's Proposition~3.8 \cite[pp.~117--118]{Bloch} says that $\lambda$ commutes with specialization; its proof permits a valuation ring with algebraically closed fraction and residue fields. Smooth proper base change gives an isomorphism on the displayed $\Q_\ell/\Z_\ell$ cohomology. Consequently the specialization of $\lambda_X^i(z)$ is $\lambda_{X_s}^i(z_s)=0$, and $\lambda_X^i(z)=0$. Geometric base change returns this conclusion to the original algebraically closed field $k$.
\end{proof}

\begin{proof}[Proof of Theorem~\ref{thm:alb}]
Put $g=\dim A$. If $u_Z=0$, every degree-zero zero-cycle on $A$ maps under $Z_*$ into $N^i(X)$. In particular
\[
 Z_*\CH_0(A)[\ell^\infty]\subseteq
                 N^i(X)\cap\CH^i(X)[\ell^\infty].
\]
Lemma~\ref{lem:torsion} and correspondence functoriality of $\lambda$ give a zero composite
\[
 \CH_0(A)[\ell^\infty]\xrightarrow{\lambda_A^g}
 H^{2g-1}(A,\Q_\ell/\Z_\ell(g))
 \xrightarrow{Z_*}H^{2i-1}(X,\Q_\ell/\Z_\ell(i)).
\]
Bloch's Proposition~3.9 and Theorem~4.2 identify the first map with the isomorphism from prime-to-$p$ torsion zero-cycles to the corresponding torsion of the Albanese variety; for $A$ this identifies it with the full displayed cohomology group. Taking $\ell$-adic Tate modules and tensoring with $\Q_\ell$ therefore gives
\[
 Z_*:H^{2g-1}(A,\Q_\ell(g))
           \longrightarrow H^{2i-1}(X,\Q_\ell(i))
\]
as a zero map. To justify this passage even when integral cohomology of $X$ has torsion, use the exact sequence from $\Z_\ell\to\Q_\ell\to\Q_\ell/\Z_\ell$: its additional term is a finite subgroup of the next integral cohomology group and has zero Tate module.

Poincar\'e duality identifies the transpose of this zero map with the map asserted in the theorem. With $\Z_\ell$ coefficients its image is torsion because its rationalization is zero. The group $H^1(A,\Z_\ell)$ is torsion free, so the integral map is zero as well.

For the last assertion take $k=\bar\F_5$ and the elliptic curve $E:y^2=x^3-x$, whose discriminant is nonzero in characteristic five. Let $A=X=E$ and $Z=\Delta_E-E\times\{0\}$. A divisor correspondence on $T\times E$ induces, after normalization at a base point, a morphism $T\to\Pic^0(E)$; this factors through $\Alb(T)$ and therefore kills its zero-cycle kernel. Hence $N^1(E)=0$, $P^1(E)=\Pic^0(E)(k)=E(k)$, and $u_Z$ is the identity. Every point is defined over a finite field, so this group is torsion and $u_Z\otimes\Q=0$. Nevertheless ${}^tZ_*$ is the identity on the nonzero two-dimensional space $H^1(E,\Q_\ell)$. Thus the integral premise cannot be replaced by its rationalization.
\end{proof}

\section{Lefschetz maps and Griffiths groups}

This section obtains the negative assertions from a single known threefold and then proves the positive and positivity statements.

\subsection{The input and the two operations}

Let $V\subset\PP^4_\C$ be a very general smooth quintic threefold. The Lefschetz theorem gives $A^1(V)=\Q h$, and $A^3(V)=\Q$ by degree. Clemens proves that the rational Abel--Jacobi image is infinite-dimensional and that algebraically trivial cycles have zero image in this case \cite[equation~(0.1), Theorem~(0.2), p.~19; genericity on p.~20]{Clemens}. Consequently
\[
 A^2(V)=\Q h^2\oplus G^2(V),\qquad
                       \dim_\Q G^2(V)=\infty.
\]
The direct sum can be chosen by the degree map; it is not canonical. All other $G^r(V)$ vanish.

Two standard formulas will be used with their functorial maps. The projective bundle formula gives
\begin{equation}\label{eq:projective}
 A^r(Y\times\PP^a)=\bigoplus_{j=0}^a A^{r-j}(Y)t^j,
 \qquad
 G^r(Y\times\PP^a)=\bigoplus_{j=0}^a G^{r-j}(Y)t^j.
\end{equation}
For a smooth codimension-two centre $Z\subset S$, with blow-up $\pi:\widetilde S\to S$ and exceptional divisor $j:E\to\widetilde S$, $q:E\to Z$, the blow-up formula is
\begin{equation}\label{eq:blowup}
 A^r(S)\oplus A^{r-1}(Z)\xrightarrow{\sim}A^r(\widetilde S),
           \quad(a,b)\longmapsto\pi^*a+j_*q^*b.
\end{equation}
The inverse is $(\pi_*,-q_*j^*)$. Indeed, $\pi_*j_*q^*=0$, $q_*j^*\pi^*=0$, and $j^*j_*q^*b=-\xi q^*b$ with $q_*\xi=1$. These operations preserve algebraic equivalence and commute with cycle classes, so \eqref{eq:blowup} also splits $G^r$. References with these maps are Jiang \cite[Theorems~2.4 and~2.8(3)]{Jiang}, and the Stacks Project \cite[Tags 02TX and 0FGQ]{Stacks}. In particular a projective bundle over a variety with a direct sum of Tate motives again has that property, and has finite-dimensional $A^r$ and zero $G^r$ in every degree.

\begin{lemma}\label{lem:induction}
For every integer $i\geq1$ there is a smooth projective $(2i+1)$-fold $X_{2i+1}$, embedded as a smooth divisor in a smooth projective $(2i+2)$-fold $R_{2i+2}$ with a direct sum of Tate motives, such that $G^r(X_{2i+1})$ vanishes except for $r=i+1$, where it is infinite-dimensional. Every other $A^r(X_{2i+1})$ is finite-dimensional.
\end{lemma}
\begin{proof}
Start with $X_3=V\subset R_4=\PP^4$. Suppose that $D=X_{2i-1}$ is a divisor in $R=R_{2i}$ with the stated properties. In
\[
 S=R\times\PP^1
\]
put $Z=D\times\{0\}$. It is the regular zero scheme of a section of
\[
 F=\operatorname{pr}_1^*\mathcal O_R(D)
          \oplus\operatorname{pr}_2^*\mathcal O_{\PP^1}(1).
\]
The blow-up $X_{2i+1}=\operatorname{Bl}_Z(S)$ is smooth. It is a divisor in the bundle of lines $R_{2i+2}=\PP_S(F)$: in local trivializations, if $(s_D,s_0)$ is the section cutting out $Z$, its equation is $s_Dv-s_0u=0$. This incidence scheme is the blow-up because the two sections form a regular sequence; it is the relative projectivization of the Rees algebra of their ideal. The equation also shows it is a Cartier divisor. The ambient bundle has a direct sum of Tate motives by the projective bundle formula.

Formula \eqref{eq:blowup} now gives
\[
 G^r(X_{2i+1})=G^r(S)\oplus G^{r-1}(Z)
                      =G^{r-1}(X_{2i-1}).
\]
Thus its only nonzero Griffiths group is in degree $i+1$, and it is infinite-dimensional. The same formula proves the finite-dimensional assertions in the other degrees. This completes the induction.
\end{proof}

\begin{proof}[Proof of Theorem~\ref{thm:range}(i) and (ii)]
For $i=0$, both $A^0(X)$ and $A^n(X)$ are one-dimensional; the map sends the fundamental class to a zero-cycle of positive degree. For $2i=n$ it is the identity.

Fix $1\leq i<n/2$ and choose $X_{2i+1}$ from Lemma~\ref{lem:induction}. Put
\[
                 Y=X_{2i+1}\times\PP^{n-2i-1}.
\]
By \eqref{eq:projective}, $A^i(Y)$ is finite-dimensional. Its $G^{n-i}(Y)$ contains the infinite-dimensional summand
\[
            G^{i+1}(X_{2i+1})t^{n-2i-1}.
\]
Hence $A^{n-i}(Y)$ is infinite-dimensional. No linear map from the former group to the latter can be surjective; this proves the assertion simultaneously for all ample classes.

In codimensions $0,1,n$, the positive assertions in part~(ii) follow respectively from connectedness, the N\'eron--Severi theorem and the degree map. For divisors, the kernel of the rational first Chern class is algebraic equivalence, and numerical equivalence gives the same rational quotient.

For every $2\leq r\leq n-1$, put $Y=V\times\PP^{n-3}$. Its $G^r(Y)$ contains $G^2(V)t^{r-2}$, an infinite-dimensional rational vector space. All those classes are numerically trivial because they are homologically trivial. Thus the cycle map is not injective, algebraic and numerical equivalence do not coincide rationally, and $A^r(Y)$ is infinite-dimensional. A finitely generated abelian group modulo torsion has finite-dimensional rationalization, so the integral finite-generation assertion fails as well.
\end{proof}

\subsection{Primitive positivity and numerical equivalence}

\begin{proof}[Proof of Theorem~\ref{thm:range}(iii)]
The cycle class of $x\in P^m_L(X)$ is a primitive Hodge class of type $(m,m)$ in $H^{2m}(X,\Q)$. The Hodge--Riemann bilinear relations say that $(-1)^m\int_X\cl(x)^2$ is positive if $\cl(x)\ne0$, and zero if $\cl(x)=0$. It follows that the form is positive semidefinite and that its radical is the kernel of the cycle map restricted to $P^m_L(X)$. This kernel is
\[
 P^m_L(X)\cap G^m(X)=\ker(L:G^m(X)\to G^{m+1}(X)).
\]
For $m=0$ the degree form is positive, and for $m=1$ the rational divisor cycle map is injective. This proves the stated positive cases.
\end{proof}

For clarity, the known low-degree homological--numerical assertion quoted in Section~1 can also be seen directly from Hodge theory. In codimension two on an $n$-fold, $n\geq4$, let $a\in H^4(X,\Q(2))$ be an algebraic class numerically trivial. Hard Lefschetz gives a unique $d\in H^2(X,\Q(1))$ with
\[
                 h^{n-2}d=h^{n-3}a.
\]
The Hodge decomposition makes $d$ of type $(1,1)$, hence a rational divisor class. Pairing this equality with divisors and using numerical triviality of $a$ shows that the Hodge-index form pairs $d$ with every rational divisor as zero. That form is nondegenerate on the divisor space, so $d=0$ and $a$ is primitive. Numerical triviality also gives $\int_Xa^2h^{n-4}=0$. The Hodge--Riemann relations then give $a=0$. The cases of curves and divisors follow in the same way from hard Lefschetz and the $(1,1)$ theorem; degrees zero and $n$ use degree. This argument uses the classical Hodge--Riemann relations and the divisor theorem; it makes no further standard-conjecture assumption.

\section{Integral scalar factorizations through Jacobians}

This section keeps track of the integer in each factorization. Losing that integer would turn the question into an isogeny statement and would not prove Theorem~\ref{thm:scalar}.

\subsection{Elementary properties and the scalar formula}

The set $I_J(A)$ is an ideal: sums use products of the intermediate products of Jacobians and integer multiples multiply the return morphism. Dualizing, using the principal polarizations of Jacobians, proves
\[
 s_J(\widehat A)=s_J(A),\qquad
 s_J(A\times B)=\operatorname{lcm}(s_J(A),s_J(B)).
\]
For the second identity, restrict and project a factorization to each factor for one inclusion, and combine two factorizations of the same integer for the other.

If $vu=1_A$ for $J=\prod J(C_j)$, the mutually inverse morphisms
\[
 A\times\ker v\longrightarrow J,\ (a,b)\longmapsto u(a)+b,
 \qquad
 J\longrightarrow A\times\ker v,\ x\longmapsto(v(x),x-uv(x))
\]
give a direct product decomposition. The kernel is proper and connected by this decomposition, and smooth because the differential of $v$ is split surjective. It is therefore an abelian variety. The converse is immediate. This proves the claimed meaning of $s_J(A)=1$.

If an isogeny $f:A\to B$ has kernel group scheme killed by $e$, the quotient property gives $h:B\to A$ with $hf=[e]_A$ and $fh=[e]_B$. Consequently
\begin{equation}\label{eq:isogeny}
 eI_J(B)\subseteq I_J(A),\qquad eI_J(A)\subseteq I_J(B).
\end{equation}
One can take $e=\deg f$, including for inseparable isogenies. The group-scheme condition, not a statement about geometric points of the kernel, is needed in positive characteristic.

\begin{lemma}[Integral Matsusaka calculation]\label{lem:scalar}
Let $X/k$ be an abelian variety of dimension $d>0$ and $L$ a line bundle with $\theta=c_1(L)$. If an integral curve cycle has rational $\ell$-adic cohomology class, for some $\ell\ne\operatorname{char}k$,
\[
                    \alpha=n\frac{\theta^{d-1}}{(d-1)!},
\]
then $n\chi(X,L)\in I_J(X)$.
\end{lemma}
\begin{proof}
Write $\alpha=\sum_jm_j(f_j)_*[C_j]$, taking smooth normalizations of its integral curve components. Choose base points and write $\bar f_j:J(C_j)\to X$ for the induced homomorphisms and $\lambda_j$ for the canonical principal polarizations. Let $\varphi_L:X\to\widehat X$ be the homomorphism associated with $L$. All morphisms in
\[
 U=(\lambda_j^{-1}\widehat{\bar f_j}\varphi_L)_j:X\to\prod_jJ(C_j),
 \qquad V=\sum_jm_j\bar f_j:\prod_jJ(C_j)\to X
\]
are integral morphisms of abelian varieties. The correspondence calculation gives
\begin{equation}\label{eq:scalar}
                         VU=[n\chi(X,L)]_X.
\end{equation}
Here is the numerical content of that calculation. A curve defines on $H^1(X,\Q_\ell)$ the alternating pairing obtained by integrating a product of two degree-one classes over the curve. Its associated symmetric map $\widehat X\to X$ is $\sum_jm_j\bar f_j\lambda_j^{-1}\widehat{\bar f_j}$. In a symplectic basis diagonalizing the alternating form of $\theta$, write
\[
             \theta=\sum_{a=1}^d t_a e_a\wedge f_a.
\]
The coefficient of the complementary pair in $\theta^{d-1}/(d-1)!$ is $\prod_{b\ne a}t_b$. Composing its map with $\varphi_L$ therefore multiplies every pair by $\prod_at_a=\chi(X,L)$, by Riemann--Roch on an abelian variety. Multiplication by $n$ proves \eqref{eq:scalar} on Tate modules. This polynomial identity is valid also when the form is degenerate. The action on a single prime-to-characteristic Tate module is faithful on homomorphisms of abelian varieties, so it proves \eqref{eq:scalar} as an equality of morphisms. This is the exact coefficient computation in the proof of Scavia--Suzuki \cite[Proposition~3.4, equations~(3.1)--(3.3)]{ScaviaSuzuki}.
\end{proof}

Apply the lemma to $X=A\times\widehat A$ and the normalized Poincar\'e line bundle $\mathcal P$. Its homomorphism exchanges the two factors, up to the conventional common sign, and is an isomorphism. Its Euler characteristic is $(-1)^g$. The integral curve cycle $c_1(\mathcal P)^{2g-1}$ corresponds to $n=(2g-1)!$. Thus $\pm(2g-1)!\in I_J(X)$; restricting to and projecting onto $A$ gives
\[
                           (2g-1)!\in I_J(A).
\]
In particular the ideal is nonzero.

\subsection{Euler characteristics and principal polarizations}

The gcd defining $\delta(A)$ is unchanged if restricted to ample line bundles. To see this, fix any $L$ and an ample $H$. The function
\[
                       p(t)=\chi(A,L\otimes H^t)
\]
is an integer-valued polynomial of degree at most $g$. The values at any $g+1$ consecutive integers generate all its integer values as an ideal: successive finite differences, or integral Newton interpolation, express every other value as an integral combination of them. Choose the consecutive integers sufficiently large that every corresponding $L\otimes H^t$ is ample. The gcd over ample Euler characteristics therefore divides $p(0)=\chi(A,L)$, which proves the assertion.

For an ample $L$ over an algebraically closed field there is a maximal isotropic finite subgroup scheme in $\ker\varphi_L$; the quotient $f:A\to B$ has degree $\chi(A,L)$ and carries a principal polarization whose pullback is $L$. This is the group-scheme form of principalization in Mumford \cite[\S23, Theorem~4 and Corollary~1, pp.~233--234]{Mumford}. It works also when the characteristic divides the degree.

When $\dim A\leq3$, the principally polarized $B$ is a product of Jacobians by Oort--Ueno \cite[Theorem~4, preprint p.~2]{OortUeno}. A compact-type curve in their statement gives exactly such a product. Equation \eqref{eq:isogeny} yields $\chi(A,L)\in I_J(A)$. Taking the gcd over ample $L$ proves $\delta(A)\in I_J(A)$.

\subsection{The factor two from generalized Pryms}

\begin{lemma}\label{lem:prym}
For a principally polarized abelian variety $(B,\Xi)$ of dimension $1\leq h\leq5$ over an algebraically closed field, there is an integral curve cycle with class
\[
                      2\frac{\xi^{h-1}}{(h-1)!},\qquad\xi=c_1(\Xi),
\]
in rational $\ell$-adic cohomology for $\ell$ different from the characteristic. In particular $2\in I_J(B)$.
\end{lemma}
\begin{proof}
For $h\leq3$, use Oort--Ueno and twice the sum of the Abel curves of the factors. Suppose $h=4$ or $5$, first in characteristic different from two. Beauville's Theorem~6.5(i) \cite[p.~178]{Beauville} realizes $(B,\Xi)$ as a generalized Prym of a nodal double cover. Let $C_a$ be the smooth components of the normalization of the covering curve. The composite from the Prym to its generalized Jacobian and then to the product of the Jacobians of these components gives
\[
                         j:B\longrightarrow J=\prod_aJ(C_a).
\]
Beauville's Proposition~2.2, p.~155, identifies the generalized theta class with the pullback of the product theta class, and his Theorem~5.4, pp.~175--176, identifies its restriction to the Prym. Thus
\[
                         j^*\Theta_J\equiv_{\alg}2\Xi.
\]
Write $\lambda_B$ and $\lambda_J$ for the principal polarizations. The polarized adjoint $q=\lambda_B^{-1}\widehat j\lambda_J$ is an integral morphism $J\to B$ and satisfies $qj=[2]_B$.

Choose Abel maps $\iota_a:C_a\to J(C_a)$ and let $q_a$ be the restriction of $q$. The integral curve $\alpha=\sum_a(q_a\iota_a)_*[C_a]$ has associated symmetric map
\[
 \sum_aq_a\lambda_a^{-1}\widehat q_a
 =q\lambda_J^{-1}\widehat q
 =\lambda_B^{-1}\widehat j\lambda_Jj\lambda_B^{-1}
 =2\lambda_B^{-1}.
\]
Pairing degree-one classes with this cycle, as in Lemma~\ref{lem:scalar}, identifies its class with $2\xi^{h-1}/(h-1)!$. Products of degree-one classes span $H^2(B,\Q_\ell)$, so those pairings determine the curve class.

In characteristic two, a principally polarized abelian variety lifts with its polarization to an abelian scheme over $W(k)$ by Oort \cite[Theorem~2.4.1 and Corollary~2.4.2, pp.~286--287]{Oort}. Only the direction for a separable polarization is needed; a principal polarization is separable. On the geometric generic fibre the preceding construction gives the integral cycle. The finite collection of its components descends to a finite extension of the fraction field. Close these components over the corresponding complete discrete valuation ring and take their special-fibre cycles. The coefficients remain integers. Specialization commutes with cycle classes for $\ell\ne2$ \cite[\S20.3, pp.~399--400]{Fulton}; hence the resulting cycle has the asserted class. Lemma~\ref{lem:scalar} with $\chi(\Xi)=1$ proves the factorization in characteristic two as well.
\end{proof}

For any abelian fourfold $X$ and every ample $L$ on it, principalization followed by Lemma~\ref{lem:prym} gives $2\chi(X,L)\in I_J(X)$. Taking the ample gcd gives $2\delta(X)\in I_J(X)$. If $A$ is a surface and $X=A\times\widehat A$, its Poincar\'e line bundle has Euler characteristic one, so $\delta(X)=1$. Hence $2\in I_J(X)\subseteq I_J(A)$. Combining this with $\delta(A)\in I_J(A)$ proves the surface bound.

\subsection{The factor six in dimension six}

\begin{lemma}\label{lem:six}
Every principally polarized abelian sixfold $(B,\Xi)$ over an algebraically closed field has an integral curve cycle with class $6\xi^5/5!$. In particular $6\in I_J(B)$.
\end{lemma}
\begin{proof}
Over $\C$ this is precisely Alexeev--Donagi--Farkas--Izadi--Ortega \cite[Corollary~2, p.~165]{ADFIO}. Their Theorem~1 on that page proves the dominance of the Prym--Tyurin construction of exponent six; the corollary is stated for every principally polarized complex sixfold.

For an arbitrary algebraically closed field $F$ of characteristic zero, define $(B,\Xi)$ over a finitely generated field $K\subset F$ and embed $K$ in $\C$. Write a complex curve cycle as $\sum_am_a(f_a)_*[C_a]$. Lemma~\ref{lem:scalar} gives the equality of homomorphisms
\begin{equation}\label{eq:sixidentity}
          \sum_am_a\bar f_a\lambda_a^{-1}\widehat{\bar f_a}\lambda_B=[6]_B.
\end{equation}
The finite collection of curves, base points, morphisms, polarizations and this equality spreads to a nonempty scheme of finite type over $K$. Shrink it to make the curves smooth projective. A closed point has residue field finite over $K$, which embeds over $K$ in $F$; its fibre gives \eqref{eq:sixidentity} over $F$. Before composition with $\lambda_B$, this identity says that the symmetric map of the curve is $6\lambda_B^{-1}$. The degree-one pairing computation therefore recovers the full class $6\xi^5/5!$, and not just an unpolarized factorization.

In positive characteristic, lift the principally polarized variety over $W(F)$ as in Lemma~\ref{lem:prym}, descend the characteristic-zero cycle to a finite extension of the fraction field, and specialize its closure. Compatibility of cycle classes gives the same equality for any $\ell$ different from the characteristic. The cycle remains integral even in characteristics two and three. Lemma~\ref{lem:scalar} then proves $6\in I_J(B)$.
\end{proof}

For an abelian sixfold $X$, principalization and the lemma give $6\chi(X,L)\in I_J(X)$ for every ample $L$, and hence $6\delta(X)\in I_J(X)$. Put $X=A\times\widehat A$ for a threefold $A$. The Poincar\'e line bundle has Euler characteristic $-1$, so $\delta(X)=1$ and $6\in I_J(A)$. Together with $\delta(A)\in I_J(A)$ this proves the threefold bound and finishes Theorem~\ref{thm:scalar}.

These gcd arguments use only finitely many factorizations. For example, on $X=A\times\widehat A$, choose an ample $H$ and $T\geq1$ so that $\mathcal P\otimes H^{T+j}$ is ample for $0\leq j\leq d=\dim X$. The explicit integer interpolation formula
\[
 \chi(\mathcal P)=\sum_{j=0}^d(-1)^j
   \binom{T+j-1}{j}\binom{T+d}{d-j}
                  \chi(\mathcal P\otimes H^{T+j})
\]
combines the available scalar factorizations into one of two or six. This is Lagrange interpolation at consecutive integers; its displayed coefficients are integers. An independent symbolic verification of the coefficient and exterior-power identities is recorded with the program \emph{Scalar identities} at \url{\workbenchurl}.

\section{Scope, related known results, and remaining questions}

Theorem~\ref{thm:alb} is one implication. It does not identify $P^i(X)$ with an ordinary or Walker intermediate Jacobian in every codimension, and it does not assert its converse. In characteristic zero, the morphic Abel--Jacobi map provides a finer obstruction: Walker's Example~5.5, Theorem~5.7, Corollary~5.9 and Theorem~6.2 \cite[pp.~928, 930, 932--934]{Walker} supply the relevant examples and functoriality. Applying that construction to Albanese-kernel-generated relations is an application of Walker's work, not a new counterexample construction.

Decomposition-of-the-diagonal arguments with zero-cycles supported on a low-dimensional subset originate in Bloch--Srinivas \cite{BlochSrinivas}. They bear directly on codimension-two Albanese and Abel--Jacobi comparisons. We do not assert a new theorem about that support condition here. In particular, the surface-supported integral comparison is not used as a premise of any proof in this paper.

For primitive positivity the exact outstanding statement is:
\[
 \ker\bigl(L:G^m(X)\longrightarrow G^{m+1}(X)\bigr)=0
 \quad\text{for every smooth projective complex }2m\text{-fold},\quad m\geq2.
\]
The examples proving Lefschetz nonsurjectivity do not settle it: the projective-space factor can make $L$ injective on the relevant Griffiths summand. The Hodge--Riemann relations control the quotient by the radical but cannot show that this radical is zero.

For Jacobian factors the unresolved cases include surfaces with even $\delta(A)$ and threefolds with $\gcd(6,\delta(A))>1$. The bounds prove neither that the residual primes occur nor that they can be removed. In complex Picard rank one, a primitive polarization of type $(1,d)$ on a surface has $\delta(A)=d$, and a type $(1,d_2,d_3)$ polarization on a threefold has $\delta(A)=d_2d_3$. Thus odd $d$, and $d_2d_3$ prime to six, give concrete positive cases. Principalization alone multiplies the scalar by the degree of its isogeny and does not resolve the residual cases.

\section*{Method and records}

The mathematics and exposition were produced by GPT-6 Astra (ChatGPT Pro mode). Separate statement-only attempts and targeted proof checks by GPT-6 Astra are recorded with their precise scope at \url{\recordurl}; the interface did not separately expose the inherited reasoning-effort setting. A statement-only attempt that did not reach a proof is recorded as such. The proofs use the named published theorems, which were checked for their stated hypotheses and locators, rather than independently reproved in full. No human has checked this paper. The original text is dedicated to the public domain under CC0 1.0. Editable sources and the statement and check records are at \url{\workbenchurl}.

\begin{thebibliography}{99}
\bibitem{Grothendieck} A. Grothendieck, \emph{Th\'eorie des motifs II}, Montpellier archive, cote~14, sheets~37--38 and~43--45. \url{https://grothendieck.umontpellier.fr/14.pdf}.
\bibitem{Beckmann} T. Beckmann and O. de Gaay Fortman, \emph{Integral Fourier transforms and the integral Hodge conjecture for one-cycles on abelian varieties}, Compos. Math. 159 (2023), 1188--1213. Author version, arXiv:2202.05230v2, Proposition~4.6, pp.~14--17, and Theorem~5.3, p.~19. \url{https://arxiv.org/abs/2202.05230}.
\bibitem{Bloch} S. Bloch, Torsion algebraic cycles and a theorem of Roitman, \emph{Compositio Math.} \textbf{39} (1979), 107--127. \url{https://www.numdam.org/article/CM_1979__39_1_107_0.pdf}.
\bibitem{Clemens} H. Clemens, Homological equivalence, modulo algebraic equivalence, is not finitely generated, \emph{Publ. Math. IH\'ES} \textbf{58} (1983), 19--38. \url{https://www.numdam.org/article/PMIHES_1983__58__19_0.pdf}.
\bibitem{Kleiman} S. L. Kleiman, Algebraic cycles and the Weil conjectures, in \emph{Dix expos\'es sur la cohomologie des sch\'emas}, North-Holland, 1968, 359--386.
\bibitem{Jiang} Q. Jiang, On the Chow theory of projectivizations, \emph{J. Inst. Math. Jussieu} \textbf{22} (2023), 1465--1508. \href{https://doi.org/10.1017/S1474748021000451}{doi:10.1017/S1474748021000451}.
\bibitem{Stacks} The Stacks Project Authors, \emph{The Stacks Project}, Tags \href{https://stacks.math.columbia.edu/tag/02TX}{02TX} and \href{https://stacks.math.columbia.edu/tag/0FGQ}{0FGQ}, accessed 23 November 2026.
\bibitem{ScaviaSuzuki} F. Scavia and F. Suzuki, \emph{On direct summands of products of Jacobians over arbitrary fields}, arXiv:2507.09821v2 (2025), especially Proposition~3.4, pp.~11--12. \url{https://arxiv.org/abs/2507.09821}.
\bibitem{Mumford} D. Mumford, \emph{Abelian Varieties}, Oxford University Press, 1970, \S23, pp.~233--234.
\bibitem{OortUeno} F. Oort and K. Ueno, Principally polarized abelian varieties of dimension two or three are Jacobian varieties, \emph{J. Fac. Sci. Univ. Tokyo IA Math.} \textbf{20} (1973), 377--381; Theorem~4 on p.~2 of the author preprint. \url{https://ir.cwi.nl/pub/27912/27912.pdf}.
\bibitem{Beauville} A. Beauville, Prym varieties and the Schottky problem, \emph{Invent. Math.} \textbf{41} (1977), 149--196. \url{https://math.univ-cotedazur.fr/u/beauvill/pubs/pv.pdf}.
\bibitem{Oort} F. Oort, Finite group schemes, local moduli for abelian varieties, and lifting problems, \emph{Compositio Math.} \textbf{23} (1971), 265--296. \url{https://www.numdam.org/item/CM_1971__23_3_265_0.pdf}.
\bibitem{Fulton} W. Fulton, \emph{Intersection Theory}, second ed., Springer, 1998, \S20.3, pp.~399--400.
\bibitem{ADFIO} V. Alexeev, R. Donagi, G. Farkas, E. Izadi and A. Ortega, The uniformization of the moduli space of principally polarized abelian 6-folds, \emph{J. reine angew. Math.} \textbf{761} (2020), 163--217. \url{https://www.mathematik.hu-berlin.de/~farkas/a6cr.pdf}.
\bibitem{Engel} P. Engel, O. de Gaay Fortman and S. Schreieder, \emph{Matroids and the integral Hodge conjecture for abelian varieties}, arXiv:2507.15704v3 (2026). \url{https://arxiv.org/html/2507.15704v3}.
\bibitem{EngelSix} P. Engel, O. de Gaay Fortman and S. Schreieder, \emph{Optimality of the Prym--Tyurin construction for $\mathcal A_6$}, preprint dated 19 December 2025, Theorem~1.2, p.~2. \href{https://www.iag.uni-hannover.de/fileadmin/iag/homepages/schreieder/publications/IHC_for_A_6.pdf}{Author manuscript}.
\bibitem{Voisin} C. Voisin, \emph{Cycle classes on abelian varieties and the geometry of the Abel--Jacobi map}, Pure Appl. Math. Q. 20 (2024), 2469--2496. Author version, arXiv:2212.03046v4, Proposition~2.1, Corollary~2.5 and Lemma~2.7, pp.~5--7. \url{https://arxiv.org/abs/2212.03046}.
\bibitem{Walker} M. E. Walker, The morphic Abel--Jacobi map, \emph{Compositio Math.} \textbf{143} (2007), 909--944. \href{https://doi.org/10.1112/S0010437X07002278}{doi:10.1112/S0010437X07002278}.
\bibitem{BlochSrinivas} S. Bloch and V. Srinivas, Remarks on correspondences and algebraic cycles, \emph{Amer. J. Math.} \textbf{105} (1983), 1235--1253. \href{https://doi.org/10.2307/2374341}{doi:10.2307/2374341}.
\bibitem{OpenAICM} OpenAI, \emph{The rational Hodge conjecture for CM abelian varieties}, preprint, 6 October 2026. \href{https://github.com/openai/math/blob/main/preprints/The-rational-Hodge-conjecture-for-CM-abelian-varieties-October-6-2026/paper.pdf}{Official preprint}, Theorem~1.1, p.~1, and Corollaries~8.3--8.4, pp.~49--50.
\bibitem{OpenAIRationality} OpenAI, \emph{Milne's rationality conjecture for abelian varieties}, preprint, 7 October 2026. \href{https://github.com/openai/math/blob/main/preprints/Milnes-rationality-conjecture-for-abelian-varieties-October-7-2026/paper.pdf}{Official preprint}, Theorem~1.1, p.~1.
\end{thebibliography}
\end{document}
